% chapters/09/evaluasi.tex
\bagian{Asesmen dan Instrumen Penilaian}
\label{sec:asesmen-fsm}

\begin{asesmenbox}
\textbf{Kedudukan asesmen bab ini.} Bab ini menuntaskan \SUBCPMK{2.2}.
Ketiga indikatornya diperiksa melalui penilaian artefak program (Tugas
Praktikum 07), diagram keadaan buatan sendiri, dan soal konsep.

\vspace{0.4em}
Materi Bab~IX termasuk bahan \textbf{Kuis 2} pada minggu kedelapan
(Materi 6--7), sebab bab ini adalah materi terakhir pada blok kuis tersebut,
dan termasuk bahan \textbf{Ujian Tengah Semester} pada minggu kesepuluh.
\end{asesmenbox}

\subsection*{Instrumen 1: Rubrik Praktikum 07}

\begin{table}[htbp]
\centering
\caption{Rubrik penilaian Praktikum 07 beserta indikator yang diukur}
\label{tab:rubrik-praktikum-07}
{\setlength{\tabcolsep}{4pt}\footnotesize%
\begin{tabularx}{\textwidth}{@{}L{3.2cm} C{1.3cm} C{2.1cm} Y@{}}
\toprule
\headrow \textbf{Kriteria} & \textbf{Bobot} & \textbf{Indikator} & \textbf{Deskripsi capaian} \\
\midrule
Diagram keadaan--transisi & 30\% & \ind{a} &
Diagram memuat empat keadaan beserta keadaan tambahan, setiap perpindahan
diberi syarat yang dapat diperiksa, dan setiap keadaan memiliki jalan keluar. \\
\addlinespace
Penerapan mesin keadaan & 35\% & \ind{b} &
Keputusan dipisahkan dari tindakan, transisi berbasis jarak dan pewaktu
diterapkan dengan benar, dan ketiga jenis pewaktu dijelaskan keperluannya. \\
\addlinespace
Pengujian dan analisis & 35\% & \ind{c} &
Catatan perpindahan beserta nomor bingkainya disajikan, skenario pengujian
dijalankan, pemeriksaan keabsahan berjalan bersih, dan hasil perpindahan yang
tidak sah dianalisis. \\
\bottomrule
\end{tabularx}}
\end{table}

\subsection*{Instrumen 2: Soal Konsep}

Jawablah dengan singkat. Kunci jawaban beserta indikator yang diukur disajikan
pada Tabel~\ref{tab:kunci-konsep-ch09}.

\begin{enumerate}
  \item Jelaskan tiga unsur yang ditambahkan mesin keadaan dibandingkan variabel
        bernilai benar-salah.
  \item Sebutkan tiga jenis pewaktu pada bab ini beserta keperluan
        masing-masing.
  \item Jelaskan mengapa fungsi \code{decide} tidak boleh mengubah keadaan apa
        pun.
  \item Sebuah NPC berada pada keadaan menyerang selama 3,6 detik dengan pewaktu
        jeda serangan 0,9 detik. Hitunglah banyaknya serangan.
  \item Mengapa pemeriksaan keabsahan perpindahan lebih berguna daripada
        pemeriksaan urutan keadaan?
\end{enumerate}

\begin{table}[htbp]
\centering
\caption{Kunci soal konsep Bab IX beserta indikator yang diukur}
\label{tab:kunci-konsep-ch09}
{\setlength{\tabcolsep}{4pt}\footnotesize%
\begin{tabularx}{\textwidth}{@{}C{0.8cm} C{2.1cm} Y@{}}
\toprule
\headrow \textbf{No} & \textbf{Indikator} & \textbf{Pokok jawaban} \\
\midrule
1 & \ind{a} &
Keadaan yang lebih dari dua beserta perilaku masing-masing; syarat perpindahan
yang tersendiri untuk setiap pasangan keadaan; dan pekerjaan penyiapan ketika
sebuah keadaan dimasuki. \\
\addlinespace
2 & \ind{b} &
Pewaktu berdiam agar keadaan berdiam tidak langsung ditinggalkan; pewaktu
tenggang agar NPC tidak langsung menyerah ketika kehilangan pemain; dan
pewaktu jeda agar serangan tidak terjadi setiap bingkai. \\
\addlinespace
3 & \ind{b} &
Agar keputusan dapat diperiksa terpisah dari tindakannya. Fungsi yang tidak
mengubah apa pun dapat dipanggil berulang kali tanpa mempengaruhi jalannya
permainan, sehingga syarat perpindahan dapat diuji langsung. \\
\addlinespace
4 & \ind{b} &
$3{,}6 \div 0{,}9 = 4$, ditambah serangan pertama yang terjadi segera setelah
keadaan dimasuki, sehingga seluruhnya 5 serangan. \\
\addlinespace
5 & \ind{c} &
Karena pemeriksaan keabsahan tidak menuntut urutan tertentu, sehingga hasilnya
tidak berubah walaupun kecepatan komputer berbeda. Pemeriksaan urutan peka
terhadap perbedaan waktu, sedangkan pemeriksaan keabsahan tetap mampu
menangkap perpindahan yang keliru. \\
\bottomrule
\end{tabularx}}
\end{table}
